\clearpage\thispagestyle{empty}
\pdfbookmark[0]{Sources for the Table for the Sun’s Coordinates (1812)}{sonnen1812-sources}
\begin{center}\textbf{Sources for the Table for the Sun’s Coordinates (1812)}\end{center}
\noindent This AI translation follows Carl Friedrich Gauss’s German text in the \href{https://gdz.sub.uni-goettingen.de/id/PPN236007467}{\textit{Werke}, vol.~VI} (Göttingen, 1874), pp.~172--180. It includes the worked example, the epochs, the Sun’s coordinates and their secular changes.
\par\smallskip\noindent The \href{https://www.digitale-sammlungen.de/en/view/bsb10538614?page=29}{first printing} appeared as article~II, \textit{Über die Tafel für die Sonnen-Coordinaten in Beziehung auf den Äquator}, in the January 1812 issue of the \textit{Monatliche Correspondenz zur Beförderung der Erd- und Himmels-Kunde}, edited by Franz Xaver von Zach, vol.~25, pp.~23--36. The \href{https://www.digitale-sammlungen.de/en/view/bsb10538614?page=34}{tables on pp.~28--36} credit \href{https://www.deutsche-biographie.de/gnd134127897.html#ndbcontent}{Friedrich Ludwig Wachter} with the calculations. Gauss’s example uses the time of an observation by von Zach and gives values independently calculated by von Lindenau.
\par\smallskip\noindent Translation and typesetting produced by AI.
\clearpage
